\chapter{Prompt de Traducción Utilizado}
\label{ap:prompt}
\begingroup
\shorthandoff{"}

Este apéndice reproduce el prompt en español utilizado en los experimentos de traducción BDD$\rightarrow$IBDD. La versión fuente empleada por el sistema puede encontrarse en el archivo \texttt{docs/PROMPT\_ES.md} del repositorio. La versión correspondiente en inglés puede encontrarse en \texttt{docs/PROMPT\_EN.md}.

\section{Contexto}

Vas a traducir escenarios escritos en formato Gherkin a un lenguaje intermedio llamado IBDD (\textit{Intermediate Behavior-Driven Development}). IBDD mantiene la estructura Given--When--Then pero añade precisión formal mediante sistemas de transición simbólicos.

\section{Definición de IBDD}

\subsection{Gramática}

\begin{quote}
\ttfamily
Scenario ::= Given When Then\\
Given ::= 'GIVEN' Declaration '[' Guard ']'\\
Declaration ::= lv1, ..., lvn\\
Guard ::= P | Guard '∧' Guard\\
When ::= 'WHEN' Switch+\\
Switch ::= Interaction Condition Assignment\\
Interaction ::= g. iv1, ..., ivn\\
Condition ::= B\\
Assignment ::= A\\
Then ::= 'THEN' Switch* '[' Guard ']'
\end{quote}

\subsection{Elementos Clave}

\begin{itemize}
\item Variables locales (\texttt{lv}): se declaran en \texttt{GIVEN} y son visibles en todo el escenario.
\item Variables globales: representan estado persistente del sistema, por ejemplo \texttt{SJL} (\textit{Scheduled Jobs List}).
\item Variables de interacción (\texttt{iv}): modelan la comunicación entre el sistema y su entorno.
\item Gates: el prefijo \texttt{!} denota salidas del sistema y el prefijo \texttt{?} denota entradas desde el entorno.
\item Guards: expresan condiciones booleanas.
\item Assignments: expresan asignaciones del tipo \texttt{variable := expresión}.
\end{itemize}

\section{Ejemplos de Traducción}

\subsection{Ejemplo 1: escenario simple con variable local}

\textbf{Escenario Gherkin}

\begin{quote}
\ttfamily\small
\{\\
\quad "id": 1,\\
\quad "domain": "printer",\\
\quad "title": "Submit job to printer",\\
\quad "given": "a job file",\\
\quad "when": "the operator submits the job file\\
\quad using Submission method",\\
\quad "then": "the printer appends a new controller job\\
\quad to the scheduled jobs AND the controller job\\
\quad is of type Job type"\\
\}
\end{quote}

\textbf{Traducción IBDD}

\begin{quote}
\ttfamily\small
GIVEN JF [true]\\
WHEN ?submit.jf,sm\\
\quad true\\
\quad JF := jf\\
THEN !append.cj,sjl\\
\quad cj.type = JT ∧ sjl = SJL ∧ JT = getJobType(SM)\\
\quad CJ := cj, SJL := add(CJ, SJL)\\
\quad [is\_in\_list(CJ, SJL) ∧ CJ.type = JT]
\end{quote}

\textbf{Razonamiento}

\begin{itemize}
\item La cláusula Given declara la variable local \texttt{JF} sin precondiciones, por lo que se utiliza la guarda \texttt{[true]}.
\item La cláusula When se modela mediante el gate de entrada \texttt{?submit} y las variables de interacción \texttt{jf} y \texttt{sm}.
\item La cláusula Then utiliza el gate de salida \texttt{!append} y una postcondición que verifica la inclusión del job en la lista.
\end{itemize}

\subsection{Ejemplo 2: variables locales con guardas y acciones secuenciales}

\textbf{Escenario Gherkin}

\begin{quote}
\ttfamily\small
\{\\
\quad "id": 2,\\
\quad "domain": "printer",\\
\quad "title": "Move production job to printed jobs",\\
\quad "given": "a controller job is in the scheduled jobs\\
\quad AND the controller job is a production job",\\
\quad "when": "the printer starts printing the controller job\\
\quad AND the printer completes printing the controller job",\\
\quad "then": "the controller job is in the printed jobs"\\
\}
\end{quote}

\textbf{Traducción IBDD}

\begin{quote}
\ttfamily\small
GIVEN CJ [is\_in\_list(CJ, SJL) ∧ CJ.type = Production]\\
WHEN !printstart.cj\\
\quad cj.id = CJ.id ∧ cj.state = Printing\\
\quad CJ.state := cj.state\\
!printcomplete.cj\\
\quad cj.id = CJ.id ∧ cj.state = Completed\\
\quad CJ.state := cj.state, PJL := add(CJ, PJL)\\
THEN [is\_in\_list(CJ, PJL)]
\end{quote}

\textbf{Razonamiento}

\begin{itemize}
\item La cláusula Given declara la variable local \texttt{CJ} y expresa una guarda compuesta sobre su pertenencia a \texttt{SJL} y su tipo.
\item La cláusula When se traduce como dos switches secuenciales, uno para el inicio de impresión y otro para su finalización.
\item La cláusula Then verifica que el job haya pasado a la lista de impresiones completadas.
\end{itemize}

\subsection{Ejemplo 3: escenario con guarda sin variables locales}

\textbf{Escenario Gherkin}

\begin{quote}
\ttfamily\small
\{\\
\quad "id": 3,\\
\quad "domain": "printer",\\
\quad "title": "Clean up printed jobs on restart",\\
\quad "given": "the printer controller is running",\\
\quad "when": "the operator restarts the printer controller",\\
\quad "then": "the printed jobs clean up is executed"\\
\}
\end{quote}

\textbf{Traducción IBDD}

\begin{quote}
\ttfamily\small
GIVEN [is\_running(CTRL)]\\
WHEN ?restart\\
\quad true\\
\quad (no assignment)\\
THEN !cleanup\\
\quad true\\
\quad PJL := empty(PJL)\\
\quad [is\_empty(PJL)]
\end{quote}

\textbf{Razonamiento}

\begin{itemize}
\item La cláusula Given no introduce variables locales; solo impone una guarda sobre el estado del controlador.
\item La cláusula When se modela mediante un gate de entrada simple.
\item La cláusula Then representa una acción del sistema seguida de una postcondición que verifica el vaciado de la lista.
\end{itemize}

\subsection{Ejemplo 4: múltiples variables de interacción}

\textbf{Escenario Gherkin}

\begin{quote}
\ttfamily\small
\{\\
\quad "id": 4,\\
\quad "domain": "ecommerce",\\
\quad "title": "Add product to cart",\\
\quad "given": "a product with price and a shopping cart",\\
\quad "when": "the user adds the product to the cart\\
\quad with quantity",\\
\quad "then": "the cart contains the product\\
\quad AND the total price is updated"\\
\}
\end{quote}

\textbf{Traducción IBDD}

\begin{quote}
\ttfamily\small
GIVEN P, SC [P.price > 0 ∧ is\_valid(SC)]\\
WHEN ?add\_to\_cart.p,sc,qty\\
\quad p.id = P.id ∧ sc.id = SC.id ∧ qty > 0\\
\quad SC.items := add(p, SC.items),\\
\quad SC.total := SC.total + (p.price * qty)\\
THEN [is\_in\_list(P, SC.items) ∧\\
\quad SC.total = calculate\_total(SC.items)]
\end{quote}

\textbf{Razonamiento}

\begin{itemize}
\item La cláusula Given introduce dos variables locales y una guarda compuesta.
\item La cláusula When utiliza múltiples variables de interacción: producto, carrito y cantidad.
\item La cláusula Then verifica tanto la inclusión del producto como la actualización correcta del total.
\end{itemize}

\section{Formato de Entrada y Salida}

La entrada del sistema consiste en un arreglo JSON de escenarios, donde cada objeto contiene los campos \texttt{id}, \texttt{domain}, \texttt{title}, \texttt{given}, \texttt{when} y \texttt{then}. La salida debe mantener exactamente esa estructura y añadir el campo \texttt{ibdd\_representation} con la traducción IBDD correspondiente.

\section{Reglas de Traducción}

\subsection{Cláusula Given}

\begin{itemize}
\item Identificar si el texto menciona variables locales, por ejemplo ``a job'' o ``a product''.
\item Si la cláusula describe únicamente el estado del sistema, modelarla solo como guarda, sin declaración.
\item Combinar múltiples condiciones mediante el operador \texttt{∧}.
\end{itemize}

\subsection{Cláusula When}

\begin{itemize}
\item Las acciones del usuario se traducen como gates de entrada.
\item Las acciones del sistema se traducen como gates de salida.
\item Las acciones múltiples conectadas por ``AND'' se traducen como switches secuenciales.
\item Deben identificarse las variables de interacción necesarias para la comunicación.
\end{itemize}

\subsection{Cláusula Then}

\begin{itemize}
\item El estado final se expresa en la postcondición.
\item Si hay acciones explícitas del sistema, deben modelarse como switches con gates de salida.
\item La postcondición final siempre se escribe entre corchetes.
\end{itemize}

\subsection{Convenciones de Nombres}

\begin{itemize}
\item Variables locales: abreviaturas en mayúsculas, por ejemplo \texttt{CJ}, \texttt{JF}, \texttt{P}, \texttt{SC}.
\item Variables globales: nombres descriptivos en mayúsculas, por ejemplo \texttt{SJL}, \texttt{PJL}, \texttt{CTRL}.
\item Variables de interacción: nombres en minúscula relacionados con las variables locales, por ejemplo \texttt{cj}, \texttt{jf}, \texttt{p}, \texttt{sc}.
\item Gates: verbos descriptivos en minúscula.
\item Ejemplos representativos son \texttt{submit}, \texttt{printstart} y \texttt{add\_to\_cart}.
\end{itemize}

\section{Instrucción Final}

Procesar el JSON de entrada y devolver el mismo JSON con el campo \texttt{ibdd\_representation} añadido a cada escenario. No deben incluirse explicaciones, comentarios ni marcas de formato adicionales; solo debe devolverse el JSON resultante, manteniendo exactamente la estructura de la entrada. La representación IBDD debe serializarse como una cadena con saltos de línea.
\endgroup
